\documentclass[10pt,a4paper]{amsart}
\usepackage{amsmath,amssymb,amsthm,graphicx}
% For commutative diagrams we use Xy-pic
\usepackage[all]{xy}

% We use 2cell for 2-commutative diagrams.
\xyoption{2cell}
% Theorem environments.
%
\theoremstyle{plain}
\newtheorem{theorem}[subsection]{Theorem}
\newtheorem{proposition}[subsection]{Proposition}
\newtheorem{lemma}[subsection]{Lemma}

\theoremstyle{definition}
\newtheorem{definition}[subsection]{Definition}
\newtheorem{example}[subsection]{Example}
\newtheorem{exercise}[subsection]{Exercise}
\newtheorem{situation}[subsection]{Situation}

\theoremstyle{remark}
\newtheorem{remark}[subsection]{Remark}
\newtheorem{remarks}[subsection]{Remarks}

\numberwithin{equation}{subsection}

% Macros
%
\def\lim{\mathop{\mathrm{lim}}\nolimits}
\def\colim{\mathop{\mathrm{colim}}\nolimits}
\def\Spec{\mathop{\mathrm{Spec}}}
\def\Hom{\mathop{\mathrm{Hom}}\nolimits}
\def\Ext{\mathop{\mathrm{Ext}}\nolimits}
\def\SheafHom{\mathop{\mathcal{H}\!\mathit{om}}\nolimits}
\def\SheafExt{\mathop{\mathcal{E}\!\mathit{xt}}\nolimits}
\def\Sch{\mathit{Sch}}
\def\Mor{\mathop{\mathrm{Mor}}\nolimits}
\def\Ob{\mathop{\mathrm{Ob}}\nolimits}
\def\Sh{\mathop{\mathit{Sh}}\nolimits}
\def\NL{\mathop{N\!L}\nolimits}
\def\CH{\mathop{\mathrm{CH}}\nolimits}
\def\proetale{{pro\text{-}\acute{e}tale}}
\def\etale{{\acute{e}tale}}
\def\QCoh{\mathit{QCoh}}
\def\Ker{\mathop{\mathrm{Ker}}}
\def\Im{\mathop{\mathrm{Im}}}
\def\Coker{\mathop{\mathrm{Coker}}}
\def\Coim{\mathop{\mathrm{Coim}}}

% Boxtimes
%
\DeclareMathSymbol{\boxtimes}{\mathbin}{AMSa}{"02}

%
% Macros for moduli stacks/spaces
%
\def\QCohstack{\mathcal{QC}\!\mathit{oh}}
\def\Cohstack{\mathcal{C}\!\mathit{oh}}
\def\Spacesstack{\mathcal{S}\!\mathit{paces}}
\def\Quotfunctor{\mathrm{Quot}}
\def\Hilbfunctor{\mathrm{Hilb}}
\def\Curvesstack{\mathcal{C}\!\mathit{urves}}
\def\Polarizedstack{\mathcal{P}\!\mathit{olarized}}
\def\Complexesstack{\mathcal{C}\!\mathit{omplexes}}
% \Pic is the operator that assigns to X its picard group, usage \Pic(X)
% \Picardstack_{X/B} denotes the Picard stack of X over B
% \Picardfunctor_{X/B} denotes the Picard functor of X over B
\def\Pic{\mathop{\mathrm{Pic}}\nolimits}
\def\Picardstack{\mathcal{P}\!\mathit{ic}}
\def\Picardfunctor{\mathrm{Pic}}
\def\Deformationcategory{\mathcal{D}\!\mathit{ef}}
\usepackage{cleveref}
\crefname{originalsection}{original section}{original sections}
\Crefname{originalsection}{Original section}{Original sections}
\makeatletter
\newlabel{section-closed-immersions}{{04E1}{source}{}{sourcepointer-1}{}}
\newlabel{section-closed-immersions@cref}{{[originalsection][0][]04E1}{[1][0][]source}{}{}{}}
\makeatother
\title{Sections with support in a closed subscheme of an injective sheaf form an injective sheaf on the closed subscheme}
\author{The Stacks Project Authors}
\date{stacks-project:etale-cohomology.tex:local-snapshot-20261011; GNU FDL 1.2 or later}
\begin{document}
\maketitle
\noindent\textit{Editorial note.} All delivered mathematical text is the author's own, copied from the pinned Stacks Project chapter file etale-cohomology.tex; no statement or proof step was rewritten, reordered or removed.\par\medskip
\noindent\textit{Editorial note.} The proof is the author's complete argument: the identification of hom groups, the exactness of i\_* and the resulting injectivity are all inside the counted span.\par\medskip
\noindent\textit{Editorial note.} The label lemma-sections-with-support-acyclic occurs in no other staged Stacks chapter, so no cross-chapter restatement had to be weighed for this card.\par\medskip
\noindent\textit{Editorial note.} The single labelled reference, which is a section pointer inside the same chapter, is carried in assembly.external\_references as a non-load-bearing navigation pointer.\par\medskip
\noindent\textit{Editorial note.} The preamble is the chapter's own native preamble excerpt already used for this source: the Xy-pic setup (native preamble lines 14-18) and the native theorem declarations with the mathematical macros (native preamble lines 154-217).\par\medskip
\noindent\textit{Allowed standard prerequisites (editorial).}\begin{itemize}
\item Sections with support in a closed subscheme. Same chapter, the sheaf H\_Z(F) on Z\_etale of sections of F supported in Z, defined by the adjunction Hom\_Z(G, H\_Z(F)) = Hom\_X(i\_*G, F). Supplies the functor whose value at an injective sheaf is examined.
\item Exactness of the direct image along a closed immersion. Same chapter, the fact that i\_* is exact for a closed immersion i : Z -> X. Supplies the exactness input that turns injectivity on X\_etale into injectivity on Z\_etale.
\item Characterisation of injective objects by exactness of the represented functor. Same chapter, the criterion that an object is injective exactly when the contravariant hom functor into it is exact. Supplies the criterion applied to H\_Z(I) on Z\_etale.
\end{itemize}\medskip
\ifdefined\hypertarget\hypertarget{sourcepointer-1}{}\fi
\noindent\textit{Source reference.} Original section 04E1: sources/stacks-etale-cohomology/pinned-original-source.tex, source lines 6575--6575. Same-chapter navigation pointer invoked at proof line 14038, where the author refers to the discussion of closed immersions for exactness of the direct image i\_*. The pointer only resolves the reference for the reader; that discussion is not reproduced inside the counted span, and no bibliography citation is involved.\par\medskip
\begin{lemma}
\label{lemma-sections-with-support-acyclic}
Let $i : Z \to X$ be a closed immersion of schemes.
Let $\mathcal{I}$ be an injective abelian sheaf on $X_\etale$.
Then $\mathcal{H}_Z(\mathcal{I})$ is an injective abelian sheaf
on $Z_\etale$.
\end{lemma}

\begin{proof}
Observe that for any abelian sheaf $\mathcal{G}$ on $Z_\etale$
we have
$$
\Hom_Z(\mathcal{G}, \mathcal{H}_Z(\mathcal{F})) =
\Hom_X(i_*\mathcal{G}, \mathcal{F})
$$
because after all any section of $i_*\mathcal{G}$ has support in $Z$.
Since $i_*$ is exact (Section \ref{section-closed-immersions}) and as
$\mathcal{I}$ is injective on $X_\etale$ we conclude that
$\mathcal{H}_Z(\mathcal{I})$ is injective on $Z_\etale$.
\end{proof}

\end{document}
